% paper/sections/07c_fccd_advection.tex
%
% §7.3  FCCD 移流：Option B/C と面ジェット共通プリミティブ
%        既存の CCD/DCCD 移流を上位互換に拡張する FCCD 移流演算子群．
%        §4.5 FCCD の面中心行列 M^FCCD を速度移流に転用する経路を示す．
%
% ═══════════════════════════════════════════════════════════════════════════════
\subsection{FCCD 移流：節点出力版 Option C と面フラックス版 Option B}
\label{sec:fccd_advection}
% ═══════════════════════════════════════════════════════════════════════════════

本節では \S\ref{sec:fccd_def} で導入した面中心 FCCD 演算子
$\mathbf{M}^{\FCCD}$（式~\eqref{eq:fccd_composite_matrix}）を
速度移流項へ拡張し，
Balanced-Force 原理（\S\ref{sec:balanced_force}）が要求する
「$\nabla p,\,\sigma\kappa\nabla\psi,\,\nabla\cdot(\mathbf{u}\otimes\mathbf{u})$
\textbf{3 項すべてが同一の面中心演算子}で離散化される」条件を完結させる．
現行実装の節点版 CCD 移流は pressure/CSF を面中心 FCCD に置換しても
\textbf{移流だけが節点中心に残る}ため BF 残差が $\mathcal{O}(h^2)$ に
退化する（本稿の数値検証による）．
これを 2 経路で解消する：
\begin{itemize}
  \item \textbf{Option C（節点出力）}：面中心 FCCD 勾配を R$_4$
        Hermite 再構成で節点に戻し，従来の AB2 履歴バッファ互換．
  \item \textbf{Option B（面フラックス）}：保存形
        $\nabla\cdot(\mathbf{u}\otimes\mathbf{u})$ を FCCD 面値 $\mathbf{P}_f$ で
        組み立て，面ダイバージェンスで節点に戻す．BF 整合性が最強．
\end{itemize}

% -------------------------------------------------------------------------------
\subsubsection{\texorpdfstring{共通プリミティブ：面値演算子 $\mathbf{P}_f$ と面ジェット $\mathcal{J}_f$}{共通プリミティブ：面値演算子 P\_f と面ジェット J\_f}}
\label{sec:fccd_adv_primitives}
% -------------------------------------------------------------------------------

節点ベクトル $\mathbf{u}\in\mathbb{R}^{N+1}$ から面ベクトル
$\mathbf{u}_f\in\mathbb{R}^N$ への 4 次精度内挿は，
節点平均の $\mathcal{O}(H^2)$ 誤差を $\mathbf{S}_{\CCD}$ が与える第二導関数
$\mathbf{q}\in\mathbb{R}^{N+1}$ で相殺することで得られる．
面中央 $x_c = x_{f_{i-1/2}}$ 周りの対称 Taylor 展開から：
\begin{equation}
  u^{\FCCD}_{f_{i-1/2}}
  = \tfrac{1}{2}(u_{i-1}+u_i)
  - \tfrac{H_i^2}{16}(q_{i-1}+q_i)
  + \mathcal{O}(H^4).
  \label{eq:fccd_adv_face_value}
\end{equation}
疎な双対角行列 $\mathbf{P}_1, \mathbf{P}_2\in\mathbb{R}^{N\times(N+1)}$ を
$\mathbf{P}_1 = \tfrac{1}{2}(1,1)\text{ 双対角}$，
$\mathbf{P}_2 = (H_i^2/16)\cdot(1,1)\text{ 双対角}$ として，
\begin{equation}
  \boxed{\;
  \mathbf{P}_f := \mathbf{P}_1 - \mathbf{P}_2\,\mathbf{S}_{\CCD},
  \qquad \mathbf{u}_f = \mathbf{P}_f\,\mathbf{u},
  \;}
  \label{eq:fccd_adv_Pf}
\end{equation}
と行列化する．
切断誤差 $u_f - u(x_c) = -\tfrac{5}{384}H^4 u^{(4)}(x_c) + \mathcal{O}(H^6)$．

\noindent\textbf{面ジェット $\mathcal{J}_f$}（\S\ref{sec:face_jet_def}）：
統一プリミティブ $\FaceJet{u} = (u_f, u'_f, u''_f)$ は
\begin{equation}
  \mathcal{J}_f(\mathbf{u}) = \bigl(\mathbf{P}_f\mathbf{u},\,
  \mathbf{M}^{\FCCD}\mathbf{u},\,
  \mathbf{S}_{\CCD}^f\mathbf{u}\bigr),
  \label{eq:fccd_adv_jet}
\end{equation}
の 3 出力を単一の $\mathbf{S}_{\CCD}$ 解から共有する．
面値・面勾配・面第二導関数はいずれも同じ $\mathbf{q}$ に依存するため
\textbf{CCD ブロック三重対角解は 1 軸あたり 1 回のみ}で済む．

% -------------------------------------------------------------------------------
\subsubsection{\texorpdfstring{Option C：R$_4$ Hermite 面→節点再構成}{Option C：R\_4 Hermite 面→節点再構成}}
\label{sec:fccd_adv_option_c}
% -------------------------------------------------------------------------------

面値勾配 $\mathbf{d}^{\FCCD} = \mathbf{M}^{\FCCD}\mathbf{u}$ から
4 次精度の節点勾配を得るには，隣接面の平均で生じる
$\mathcal{O}(H^2)$ 誤差を $\mathbf{q}$ で相殺する：
\begin{equation}
  (\partial_x u)^{\text{node-FCCD}}_i
  = \tfrac{1}{2}(d_{f_{i-1/2}} + d_{f_{i+1/2}})
  - \tfrac{H}{16}(q_{i+1} - q_{i-1})
  + \mathcal{O}(H^4).
  \label{eq:fccd_adv_R4}
\end{equation}
この\textbf{R$_4$ Hermite 再構成}で節点勾配が FCCD の面精度と
同じ $\mathcal{O}(H^4)$ を達成する．
移流項は従来通り node-pointwise 積で組み立てる：
\begin{equation}
  [\mathcal{A}^{(j)}]_i = -\sum_k u_k(x_i)\cdot(\partial_{x_k} u_j)^{\text{node-FCCD}}_i.
  \label{eq:fccd_adv_option_c}
\end{equation}

\noindent\textbf{Option C の利点}：
出力が節点値なので既存の AB2 履歴バッファ / PPE RHS /
CSF $\sigma\kappa\nabla\psi$ / Rhie--Chow がすべて無変更．
`ns\_pipeline.py` の \texttt{ccd.differentiate(u, ax)[0]} を
\texttt{fccd.node\_gradient(u, ax, q=...)} に置換するだけで実装完了．

\noindent\textbf{Option C の限界}：
R$_4$ の面→節点変換で BF 整合性が\textbf{部分的に}回復するに留まる．
圧力と CSF が面中心 FCCD で動く一方，移流だけは R$_4$ 経由の
節点に戻っているため，節点-面-節点の往復で
$\mathcal{O}(H^4)$ 誤差が 2 重に蓄積する．
製品実装では Option B が推奨．

% -------------------------------------------------------------------------------
\subsubsection{Option B：保存形 FCCD 面フラックス}
\label{sec:fccd_adv_option_b}
% -------------------------------------------------------------------------------

運動量の自然形は $\nabla\cdot(\mathbf{u}\otimes\mathbf{u})$（保存形）である．
非保存形 $(\mathbf{u}\cdot\nabla)\mathbf{u}$ は $\nabla\cdot\mathbf{u}\neq 0$
の離散的残差によりスプリアス運動量源を生む．
FCCD 面フラックスで保存形を評価する：
\begin{equation}
  F^{(k,j),\text{cons}}_{f_{i-1/2}}
  = (u^{(k)} u^{(j)})_f
  = \bigl(\mathbf{P}_f(u^{(k)} u^{(j)})\bigr)_{i-1/2}
  \quad\text{（節点積 $u^{(k)}_i u^{(j)}_i$ を先に作り，単一の $\mathbf{P}_f$ で面に補間）},
  \label{eq:fccd_adv_flux_cons}
\end{equation}
面で 2 回分離補間する $(\mathbf{P}_f u^{(k)})(\mathbf{P}_f u^{(j)})$ は
非線形積の補間誤差が $\Ord{H^2}$ フロアに落ちるため採用しない．
節点ダイバージェンスで運動量更新を閉じる：
\begin{equation}
  [\nabla\cdot F^{(j)}]_i
  = \sum_k\frac{F^{(k,j)}_{f_{i+1/2}} - F^{(k,j)}_{f_{i-1/2}}}{H_i}.
  \label{eq:fccd_adv_divergence}
\end{equation}

\noindent\textbf{スキュー対称形（推奨）}：
運動量の離散エネルギー保存のため，保存形と非保存形の算術平均
\begin{equation}
  F^{(k,j),\text{sk}}_{f_{i-1/2}}
  = \tfrac{1}{2}\bigl(F^{(k,j),\text{cons}}_{f_{i-1/2}}
  + F^{(k,j),\text{nc}}_{f_{i-1/2}}\bigr)
  \label{eq:fccd_adv_flux_skew}
\end{equation}
を用いる（Kok 2000, Ham et al. 2002）．
非保存形 $F^{(k,j),\text{nc}} = (\mathbf{P}_f u^{(k)})\cdot(\mathbf{M}^{\FCCD}u^{(j)})$
は面値と面勾配の点積，両者とも $\mathcal{J}_f$ から得られる．

% -------------------------------------------------------------------------------
\subsubsection{BF 整合性定理（面上完結）}
\label{sec:fccd_adv_bf_theorem}
% -------------------------------------------------------------------------------

\begin{theorem}[FCCD Option B の BF 整合性]\label{thm:fccd_adv_bf}
$\nabla p$（FCCD），$\sigma\kappa\nabla\psi$（FCCD 面形），
$\nabla\cdot(\mathbf{u}\otimes\mathbf{u})$（FCCD Option B）の
3 項すべてを\textbf{共通の面中心演算子}
$\mathbf{M}^{\FCCD}$（\S\ref{sec:fccd_def} 式~\eqref{eq:fccd_composite_matrix}）で
離散化すれば，静止状態 $\mathbf{u}\equiv 0$ での離散運動量バランスの
BF 残差は
\begin{equation*}
  \|\mathbf{R}_\text{BF}\|_\infty = \mathcal{O}(H^4)\text{（一様）},
  \qquad \mathcal{O}(H^3)\text{（非一様）}
\end{equation*}
を満たす．
\end{theorem}

\begin{proof}[証明（要旨）]
本定理は PPE が $[p]=\sigma\kappa$ を\textbf{有限誤差 $\Ord{H^4}$ 内で解く}
（すなわち PPE 残差 $\|G_h p - f_\sigma\|_\infty = \Ord{H^4}$；
\S\ref{sec:split_ppe_pressure_jump_formulation}）ことを仮定する．
$\mathbf{u}\equiv 0$ で保存形・非保存形・スキュー対称形の 3 種の対流面
フラックス（\eqref{eq:fccd_adv_flux_cons}
\eqref{eq:fccd_adv_flux_skew} および非保存形）はいずれも
点積に $\mathbf{u}$ を含むため恒等的に零．
残る $\nabla p$ と CSF は共通演算子 $\mathbf{M}^{\FCCD}$ の
切断誤差 $\mathcal{O}(H^4)/\mathcal{O}(H^3)$
（\S\ref{sec:fccd_def}/\S\ref{sec:fccd_nonuniform}）で相殺する．$\square$
\end{proof}

\noindent\textbf{系}：
移流を Option C（R$_4$ 経由の節点中心）に留めた場合，
節点-面-節点の往復で BF 残差フロアは $\mathcal{O}(H^2)$ に退化．
Option B でのみ H-01 が完全解消する（本稿の実測で移流のみ節点中心 CCD の
$\mathcal{O}(H^2)$ フロアから $\mathcal{O}(H^4)$ へ降下）．

% -------------------------------------------------------------------------------
\subsubsection{運動量方程式の閉形}
\label{sec:fccd_adv_closure}
% -------------------------------------------------------------------------------

FCCD Option B で閉じた運動量更新：
\begin{equation}
  \frac{u^{n+1}_j - u^n_j}{\Delta t}
  = -\sum_k\frac{(u_k u_j)^{*}_{f_{i+1/2}} - (u_k u_j)^{*}_{f_{i-1/2}}}{H_i}
  + \bigl[\nu \text{Lap} + \nabla p + \sigma\kappa\nabla\psi + g\bigr]_i,
  \label{eq:fccd_adv_closure}
\end{equation}
右辺は粘性・圧力・CSF・重力すべて\textbf{面中心}で評価し
node\_divergence で節点に戻す．
PPE RHS は自然に $\nabla_h\cdot\mathbf{u}^*$ が FCCD 面ダイバージェンスで
表現され，ループが閉じる．

% -------------------------------------------------------------------------------
\subsubsection{壁境界条件 Option IV：Dirichlet 無滑り}
\label{sec:fccd_adv_wall_iv}
% -------------------------------------------------------------------------------

\S\ref{sec:ccd_bc} で導入した Option III（Neumann $\psi, p$）は
移流でも変更不要．
Dirichlet 無滑り $u(x_\text{wall}) = 0$ に対する
\textbf{Option IV（Dirichlet 壁境界）}は，
$u$ の反対称ゴーストミラー $u_{-1} = -u_1$ と
$u''$ の反対称性 $q_{-1} = -q_1$ を組み合わせる：
\begin{align}
  d^{\FCCD}_{f_{-1/2}} &= \frac{u_1}{H} - \frac{H}{24}(q_0 + q_1)
  \quad\text{（壁せん断率, 一般に非零）}, \label{eq:fccd_adv_wall_iv_d}\\
  u^{\FCCD}_{f_{-1/2}} &= -\tfrac{1}{2}u_1 + \tfrac{H^2}{16}(q_1 - q_0)
  \quad\text{（壁面速度, 連続無滑り極限で零）}. \label{eq:fccd_adv_wall_iv_u}
\end{align}
Option B の壁面フラックス
$F^{(k,j)}_{f_{-1/2}} = u^{(k)}_f\cdot u^{(j)}_f$ は
両因子とも壁で $\to 0$ のため先頭オーダーで恒等零，
物理的無滑り運動量フラックス零と一致．

% -------------------------------------------------------------------------------
\subsubsection{実装・コスト・段階移行}
\label{sec:fccd_adv_impl}
% -------------------------------------------------------------------------------

\noindent\textbf{$\mathbf{q}$-共有最適化}：
面値 $\mathbf{P}_f\mathbf{u}$，面勾配 $\mathbf{M}^{\FCCD}\mathbf{u}$，
面第二導関数 $\mathbf{S}_{\CCD}^f\mathbf{u}$ はすべて共通の
$\mathbf{q} = \mathbf{S}_{\CCD}\mathbf{u}$ に依存するため，
CCD ブロック三重対角解は\textbf{1 軸あたり 1 回のみ}．
追加コストは疎行列 $\mathbf{P}_1, \mathbf{P}_2, \mathbf{D}_1, \mathbf{D}_2$ の
$\mathcal{O}(N)$ 乗算のみで，
演算量は純 CCD 移流の 1.3--1.5$\times$．

\noindent\textbf{段階移行パス}：
\begin{enumerate}
  \item \textbf{Stage 1}：pressure/CSF のみ FCCD 化（\S\ref{sec:balanced_force}），
        移流は従来の節点 CCD を維持．
        BF 残差は pressure/CSF 精度に律速され $\mathcal{O}(H^2)$ 止まり．
  \item \textbf{Stage 2}：Option C 導入で移流を R$_4$ 経由の節点中心 FCCD に．
        $\mathcal{O}(H^2)$ フロアは残るが係数が小さくなる．
  \item \textbf{Stage 3}：Option B で保存形面フラックスに完全移行．
        H-01 完全解消，BF 残差 $\mathcal{O}(H^4)$．
\end{enumerate}
製品コンフィグは Stage 3，
回帰検証には Stage 1/2 を \texttt{fccd\_bf\_mode} フラグで切替可能．

% ═══════════════════════════════════════════════════════════════════════════════
